%------------------------------
% Resume in Latex
% Author : Zijie Zhang
%------------------------------

%------------------------------
% DOCUMENT DEFINITION
%------------------------------
\documentclass[UTF8,A4]{ctexart}


%------------------------------
% ENCODING
%------------------------------
%we use utf8 since we want to build from any machine
\usepackage[utf8]{inputenc}


%------------------------------
% FONT
%------------------------------
% \usepackage{helvet}
% \usepackage[default]{raleway}

% set font default
\renewcommand*\familydefault{\sfdefault} 	
\usepackage[T1]{fontenc}

% more font size definitions
\usepackage{moresize}


%------------------------------
% PAGE LAYOUT DEFINITIONS
%------------------------------
\usepackage[a4paper]{geometry}
\geometry{top=1.25cm, bottom=-.6cm, left=1.5cm, right=1.5cm}

%use customized header
\usepackage{fancyhdr}
\pagestyle{fancy}

%less space between header and content
\setlength{\headheight}{-5pt}

%indentation is zero
\setlength{\parindent}{0mm}


%------------------------------
% TABLE/ARRAY DEFINITIONS
%------------------------------
%for layouting tables
\usepackage{multicol}
\usepackage{multirow}

%extended aligning of tabular cells
\usepackage{array}

\newcolumntype{x}[1]{%
>{\raggedleft\hspace{0pt}}p{#1}}%


%------------------------------
% GRAPHICS DEFINITIONS
%------------------------------
%for header image
\usepackage{graphicx}

%for floating figures
\usepackage{wrapfig}
\usepackage{float}

%for drawing graphics
\usepackage{tikz}
\usetikzlibrary{shapes, backgrounds,mindmap, trees}


%------------------------------
% COLOR DEFINITIONS
%------------------------------
\usepackage{color}

%accent color
\definecolor{sectcol}{RGB}{0,150,255}

%dark background color
\definecolor{bgcol}{RGB}{110,110,110}

%light background / accent color
\definecolor{softcol}{RGB}{225,225,225}



%==================================================
% DEFINITIONS
%==================================================

%-------------------------
% HEADER
%-------------------------
% remove top header line
\renewcommand{\headrulewidth}{0pt}

%remove botttom header line
\renewcommand{\footrulewidth}{0pt}

%remove pagenum
\renewcommand{\thepage}{}

%remove section num
\renewcommand{\thesection}{}


%-------------------------
% ARROW GRAPHICS in Tikz
%-------------------------
% a six pointed arrow poiting to the left
\newcommand{\tzlarrow}
    {(0,0) -- (0.2,0) -- (0.3,0.2) -- (0.2,0.4) -- (0,0.4) -- (0.1,0.2) -- cycle;}	

% include the left arrow into a tikz picture
% param1: fill color
%
\newcommand{\larrow}[1]
{
\begin{tikzpicture}[scale=0.58]
    \filldraw[fill=#1!100,draw=#1!100!black]  \tzlarrow
\end{tikzpicture}
}

% a six pointed arrow poiting to the right
\newcommand{\tzrarrow}
    { (0,0.2) -- (0.1,0) -- (0.3,0) -- (0.2,0.2) -- (0.3,0.4) -- (0.1,0.4) -- cycle;}

% include the right arrow into a tikz picture
% param1: fill color
%
\newcommand{\rarrow}
{
\begin{tikzpicture}[scale=0.7]
	\filldraw[fill=sectcol!100,draw=sectcol!100!black] \tzrarrow
\end{tikzpicture}
}



%-------------------------
% CUSTOM SECTIONS
%-------------------------
% create a coloured box with arrow and title as cv section headline
% param 1: section title
%
\newcommand{\cvsection}[1]
{
	\begin{center}
		\large\textcolor{sectcol}{\textbf{#1}}
	\end{center}
}

%create a coloured arrow with title as cv meta section section
% param 1: meta section title
%
\newcommand{\metasection}[2]
{
%\begin{tabular*}{1\textwidth}{r r}
\footnotesize{#2} \hspace*{\fill} \footnotesize{#1}\\[1pt]
%\end{tabular*}
}


%-------------------------
% CV EVENT
%-------------------------

\newcommand{\cvevent}[3]
{

\begin{tabular*}{1\textwidth}{p{13.6cm}  x{3.9cm}}
	\textbf{#2} - \textcolor{bgcol}{#3} &   \vspace{2.5pt}\textcolor{sectcol}{#1}
\end{tabular*}

\vspace{-8pt}
\textcolor{softcol}{\hrule}
\vspace{6pt}

    % $\cdot$ #4\\[3pt]
    % $\cdot$ #5\\[6pt]

}

% creates a stretched box as 
\newcommand{\cveventmeta}[2]
{
	\mbox{\mystrut \hspace{87pt}\textit{#1}}\\
	#2
}


%-------------------------
% CUSTOM STRUT FOR EMPTY BOXES
%-------------------------
\newcommand{\mystrut}{\rule[-.3\baselineskip]{0pt}{\baselineskip}}



%==================================================
% DOCUMENT CONTENT
%==================================================

\begin{document}


%use our custom fancy header definitions
\pagestyle{fancy}

%-------------------------
% TITLE HEADLINE
%-------------------------
\vspace{-8pt}
\begin{center}
    \HUGE   \textsc{张子杰}
            \textcolor{sectcol}{\rule[-1mm]{1mm}{0.9cm}}
            \textsc{简历}\\[2pt]
	\small  意向部门: \textit{技术研究，机器学习方向}
\end{center}

\vspace{6pt}

%-------------------------
% META SECTION
%-------------------------
\metasection{Madison, WI, United States}{\textbf{Status:} M.A. in Mathematics, Postgraduate}
\metasection{github.com/z-zijie}{\textbf{Fields:} 机器学习, 计算数学, 数据科学} 
\metasection{z.zijie@outlook.com}{}

\vspace{-2pt}
\textcolor{softcol}{\hrule}
\vspace{6pt}

\normalsize



%==================================================
% CV SECTIONS AND EVENTS (MAIN CONTENT)
%==================================================

%-------------------------
% EDUCATION SECTION
%-------------------------
\cvsection{EDUCATION}

\cvevent{9月 2019 - 5月 2021}
        {M.A. in 数学硕士}
        {威斯康辛大学麦迪逊分校}

        $\cdot$ 主要学习方向 计算数学，数据科学，机器学习和随机分析\\
        $\cdot$ GPA: 3.88/4.0 \\

\cvevent{9月 2019 - 5月 2020}
        {出国交换}
        {威斯康辛大学麦迪逊分校}
        $\cdot$ 主要学习方向 应用数学与统计学习\\
        $\cdot$ GPA: 3.375/4.0 \\

\cvevent{9月 2016 - 5月 2020}
        {统计学学士}
        {南京大学}

        $\cdot$ GPA.Cumulative: 3.95/5
        $\cdot$ GPA.Major: 3.91/5\\


%-------------------------
% RESEARCH SECTION
%-------------------------
\cvsection{RESEARCH}

\cvevent{9月 2018 – 12月 2018}
        {南京大学金融创新实验室}
        {南京, 江苏, 中国}

        $\cdot$ 这是关于区块链的算法设计科研经历。\\
        $\cdot$ 作为团队的一员，为高速交易设计了一个共轭routing和账户平衡算法，并撰写了论文加以阐述。\\
        $\cdot$ 在天河超算上用MATLAB进行仿真和数值优化，使交易成功率超过90％。\\
        $\cdot$ 该项目的论文（三作）， 
                "Photon State-Channel Architecture with AI Routing Optimization", 
                已被2018 CCF Blockchain conference接受。\\

\cvevent{Feb 2018}
        {美国大学生数学建模竞赛（交叉学科）}
        {南京, 江苏, 中国}

        $\cdot$ 设计了一个模型来确定特斯拉目的地充电站和超级充电站的数量、分布和时间表。\\
        $\cdot$ 使用排队模型来计算预期等待时间和到达率之间的关系。\\
        $\cdot$ 采用线性回归模型对电动汽车的数量进行预测。\\
        $\cdot$ 该项目产生的论文被授予Honorable Mention in ICM.\\



%-------------------------
% EXPERIENCE SECTION
%-------------------------
\cvsection{EXPERIENCE}
\cvevent{NOV 2020}
        {口译}
        {威斯康辛大学麦迪逊分校}

        $\cdot$ 为南京大学代表团提供口译服务。\\
        $\cdot$ 为与Exact Sciences生物医药癌症检测技术交流会参会者提供实时口译。\\
        $\cdot$ 陪同参观并即时翻译交流，位于威斯康星州麦迪逊市的研究所，Discovery Building.\\


\newpage


%-------------------------
% COURSEWORK HIGHLIGHTS
%-------------------------
\cvsection{COURSEWORK HIGHLIGHTS}
\vspace{-18pt}
\textcolor{softcol}{\hrule}
\vspace{6pt}

\textbf{数学}
实分析 \textbullet{} 复分析 \textbullet{} 泛函分析 \\
应用数学 I \& II \textbullet{} 计算数学 I \& II \\
运筹学 \textbullet{} 抽象代数 \textbullet{} 应用数学主题\\

\textbf{统计学}
统计学习 \textbullet{} 随机过程 \textbullet{} 数理统计 \\
多元统计分析 \textbullet{} 时间序列分析 \textbullet{} 风险统计\\

\textbf{计算机科学}
模式识别 \textbullet{} 数据挖掘 \textbullet{} 机器学习\\


%-------------------------
% SKILLS
%-------------------------
\cvsection{SKILLS}
\vspace{-18pt}
\textcolor{softcol}{\hrule}
\vspace{6pt}

\textbf{PROGRAMMING}\\
Python \textbullet{} Matlab \textbullet{} C \textbullet{} C++ \textbullet{} \LaTeX\ \\
R \textbullet{} Shell \textbullet{} Fortran(w/ mpi) \textbullet{} Octave\\
\textit{Familiar:}
Java \textbullet{} Assembly \textbullet{} MySQL\\

\textbf{SPOKEN \& WRITTEN}\\
流利: 英语 \\
母语: 中文\\

\end{document}